% =====================================================================
% Capítulo 4 · Análisis
% Borrador generado a partir de docs/tfg/REQUISITOS.md, TRAZABILIDAD.md,
% DECISIONES.md (ADR-004, ADR-005, ADR-011, ADR-014) y AGENTS.md.
% Estado de las fuentes: 2026-09-29 (F2.4b cerrada).
% Paquetes necesarios: booktabs, tabularx, todonotes, tikz
%   (bibliotecas: positioning, shapes.geometric, fit, backgrounds).
% Alternativa en PlantUML: diagramas/casos-uso.puml
% =====================================================================

\chapter{Análisis}
\label{cap:analisis}

Este capítulo delimita el problema que resuelve PymeKit, recoge sus
requisitos funcionales y no funcionales y describe los casos de uso con los
que se validará la plataforma.

% ---------------------------------------------------------------------
\section{Contexto del problema}
\label{sec:contexto-problema}

Las pequeñas y medianas empresas (pymes) consumen cada vez más aplicaciones en
modalidad \textit{Software as a Service} (SaaS, software como servicio): el
proveedor aloja y mantiene una única instalación y cada cliente accede a ella
mediante una suscripción. Quien desarrolla un SaaS para pymes debe resolver,
antes de escribir una sola línea de su lógica de negocio, un conjunto de
problemas que se repiten en casi todos los productos de este tipo:

\begin{itemize}
  \item \textbf{identidad}: registro, inicio de sesión, recuperación de
        contraseña y, cada vez más a menudo, un segundo factor de
        autenticación;
  \item \textbf{multi-tenencia} (\textit{multi-tenancy}): cada pyme cliente
        es un \textit{tenant} (inquilino) cuyos datos deben quedar aislados de
        los de los demás, aunque todos compartan la misma base de datos;
  \item \textbf{organización interna del cliente}: equipos, invitaciones,
        roles y permisos;
  \item \textbf{cobro}: planes, suscripciones y sincronización con una
        pasarela de pago;
  \item \textbf{operación de la plataforma}: una consola para que el
        personal del proveedor gestione usuarios y cuentas, y una herramienta
        para consultar y corregir datos sin acceder directamente a la base de
        datos;
  \item \textbf{calidad transversal}: seguridad, internacionalización,
        pruebas y despliegue reproducible.
\end{itemize}

Estos componentes no aportan valor diferencial al producto, pero consumen una
parte muy importante del esfuerzo y son precisamente los que más riesgo de
seguridad concentran. El objetivo de PymeKit es ofrecerlos ya resueltos,
integrados y verificados, como una \textbf{plataforma base reutilizable}: un
nuevo SaaS para pymes debería poder arrancarse cambiando la configuración, sin
modificar el núcleo (RNF-01).

Del alcance acordado se derivan dos decisiones que condicionan el análisis:

\begin{itemize}
  \item La plataforma es \textbf{genérica}: no implementa un dominio de
        negocio concreto. Su reutilización se valida midiendo el arranque de
        un SaaS nuevo a partir de ella (ADR-005).
        % PENDIENTE (P-01): confirmar con el tutor que esto cubre el
        % «escenario de ejemplo» de la propuesta.
  \item El \textbf{CMS} que pide la propuesta se interpreta como un CMS de
        administración de datos de la base de datos, integrado en la consola
        de administración, y no como un gestor de contenidos editoriales
        (blog, documentación). Mantener dos CMS añadiría complejidad sin
        acercar el objetivo (ADR-004, ADR-011).
\end{itemize}

% PENDIENTE (P-06): si la memoria explica aquí que la plataforma parte de una
% base de código SaaS existente con licencia, redactarlo de forma neutra (sin
% nombres comerciales ni URL) según lo que se acuerde con el tutor. La parte
% reutilizada y la propia se distinguen en el capítulo de Implementación.

% ---------------------------------------------------------------------
\section{Actores}
\label{sec:actores}

\begin{description}
  \item[Visitante.] Persona no identificada que consulta la web pública y
        puede registrarse.
  \item[Usuario.] Persona registrada. Dispone de una cuenta personal y puede
        pertenecer a una o varias cuentas de equipo.
  \item[Propietario del equipo.] Usuario con el rol de mayor nivel en una
        cuenta de equipo (la pyme cliente); gestiona sus miembros y su
        suscripción.
  \item[Miembro.] Usuario de un equipo con un rol de menor nivel y, por
        tanto, con permisos limitados.
  \item[Super-administrador.] Personal del proveedor con control total de la
        plataforma y, por extensión, del CMS (ADR-014).
  \item[Gestor de contenidos / soporte.] Personal del proveedor que accede al
        CMS con permisos limitados definidos por el RBAC propio del CMS
        (ADR-014, ADR-016).
  \item[Desarrollador.] Quien reutiliza PymeKit para construir un SaaS nuevo.
  \item[Pasarela de pago.] Sistema externo (Stripe) que procesa los cobros y
        notifica los cambios de suscripción mediante \textit{webhooks}
        (llamadas HTTP que el sistema externo envía a la aplicación cuando
        ocurre un evento).
\end{description}

% ---------------------------------------------------------------------
\section{Requisitos funcionales}
\label{sec:rf}

Los requisitos se derivan de los apartados \emph{Objetivo}, \emph{Descripción}
y \emph{Alcance} de la propuesta del TFG. Su prioridad es \textbf{M}
(imprescindible, forma parte del prototipo mínimo) o \textbf{D} (deseable).
Los identificadores son estables: se citan en el código, en la matriz de
trazabilidad y en esta memoria, y no se renumeran.

\begin{table}[htbp]
  \centering
  \footnotesize
  \caption{Requisitos funcionales.}
  \label{tab:rf}
  \begin{tabularx}{\textwidth}{@{}lXc@{}}
    \toprule
    \textbf{ID} & \textbf{Requisito} & \textbf{Prior.} \\
    \midrule
    RF-01 & Interfaz web para el usuario final: página de inicio pública (\textit{landing}), página de precios y área privada con navegación & M \\
    RF-02 & API de \textit{backend} para la lógica de negocio y el acceso a datos, mediante funciones de servidor (\textit{server functions}) tipadas y validadas & M \\
    RF-03 & Registro e inicio de sesión (contraseña y enlace mágico), recuperación de contraseña y verificación del correo electrónico & M \\
    RF-04 & Segundo factor de autenticación (MFA) y proveedores OAuth configurables & D \\
    RF-05 & Gestión de la cuenta personal: perfil, correo, contraseña, preferencia de idioma y baja & M \\
    RF-06 & Cuentas de equipo (una por pyme cliente): crear equipos, invitar miembros, asignar roles y aplicar permisos por funcionalidad & M \\
    RF-07 & Suscripciones con pasarela de pago (Stripe): planes, pago (\textit{checkout}), portal de cliente y sincronización mediante \textit{webhooks} & M \\
    RF-08 & Panel de super-administración de la plataforma: listar y gestionar usuarios y cuentas, bloquear, suplantar y restablecer & M \\
    RF-09 & CMS de administración de datos, integrado en la consola de administración: explorar y editar tablas de la BD, gestionar usuarios y almacenamiento, con permisos propios & M \\
    RF-10 & Registro de auditoría de las acciones realizadas desde el CMS & D \\
    RF-11 & Paneles (\textit{dashboards}) configurables en el CMS & D \\
    RF-12 & Notificaciones dentro de la aplicación y correos transaccionales & D \\
    RF-13 & Interfaz multidioma: español por defecto e inglés & M \\
    \bottomrule
  \end{tabularx}
\end{table}

RF-09 es el requisito que más ha evolucionado durante el análisis: en un
principio se planteó el CMS como una aplicación separada que compartía la
base de datos (ADR-002) y, tras evaluar alternativas, se decidió integrarlo
en la consola de administración (ADR-011). De ahí que su redacción incluya
«integrado en la consola de administración». RF-11 es el primer candidato a
recortar si el plazo lo exige.

% ---------------------------------------------------------------------
\section{Requisitos no funcionales}
\label{sec:rnf}

Cada requisito no funcional se acompaña del modo en que se verifica, para que
sea comprobable y no una mera declaración de intenciones.

\begin{table}[htbp]
  \centering
  \footnotesize
  \caption{Requisitos no funcionales y forma de verificarlos.}
  \label{tab:rnf}
  \begin{tabularx}{\textwidth}{@{}l>{\raggedright\arraybackslash}p{6.2cm}X@{}}
    \toprule
    \textbf{ID} & \textbf{Requisito} & \textbf{Verificación} \\
    \midrule
    RNF-01 & \textbf{Reutilización y configurabilidad}: la plataforma se adapta a un nuevo SaaS cambiando la configuración (ficheros de configuración y variables de entorno), sin tocar el núcleo & Evaluación de la reutilización (F6): pasos y tiempo para arrancar un proyecto nuevo \\
    RNF-02 & \textbf{Seguridad}: la autorización se aplica en la base de datos (RLS) además de en la aplicación; secretos fuera del repositorio y validación de todas las entradas & Pruebas pgTAP y revisiones adversariales de RLS y del código \\
    RNF-03 & \textbf{Aislamiento multi-tenant}: los datos de una cuenta no son accesibles desde otra & Pruebas pgTAP específicas de aislamiento \\
    RNF-04 & \textbf{Modularidad}: monorepo de paquetes desacoplados, con dependencias explícitas y exportaciones separadas para cliente y servidor & Estructura del monorepo, comprobación de dependencias y de tipos \\
    RNF-05 & \textbf{Mantenibilidad}: TypeScript estricto, \textit{lint} y formato automáticos, y código comentado en español didáctico & Integración continua y revisión de comentarios \\
    RNF-06 & \textbf{Calidad verificable}: pruebas unitarias, de base de datos y E2E de los casos de uso principales & Integración continua en verde \\
    RNF-07 & \textbf{Desplegabilidad}: despliegue reproducible con contenedores y en la nube, con manual de instalación & Despliegue de prueba (F7) \\
    RNF-08 & \textbf{Usabilidad y accesibilidad}: interfaz adaptable (\textit{responsive}), componentes accesibles y tema claro/oscuro & Revisión manual y E2E \\
    RNF-09 & \textbf{Rendimiento razonable}: carga de datos en el servidor, caché de consultas y paginación en el CMS & Revisión manual \\
    \bottomrule
  \end{tabularx}
\end{table}

La relación entre cada requisito, el código que lo implementa y las pruebas
que lo verifican se mantiene en una matriz de trazabilidad que se actualiza al
cerrar cada incremento.
\todo{Incluir la matriz de trazabilidad (requisito $\rightarrow$ módulo
$\rightarrow$ pruebas $\rightarrow$ estado) como anexo o en el capítulo de
Pruebas, con el estado final tras F6.}

% ---------------------------------------------------------------------
\section{Casos de uso}
\label{sec:casos-uso}

Se han seleccionado siete casos de uso principales que recorren los módulos
de la plataforma y que servirán de base para la validación de la fase F6
(tabla~\ref{tab:cu} y figura~\ref{fig:casos-uso}).

\begin{table}[htbp]
  \centering
  \footnotesize
  \caption{Casos de uso principales.}
  \label{tab:cu}
  \begin{tabularx}{\textwidth}{@{}l>{\raggedright\arraybackslash}p{2.6cm}Xl@{}}
    \toprule
    \textbf{CU} & \textbf{Actor} & \textbf{Descripción} & \textbf{Requisitos} \\
    \midrule
    CU-01 & Visitante & Se registra, verifica el correo y accede al área privada & RF-01, RF-03 \\
    CU-02 & Usuario & Crea un equipo para su pyme e invita a un empleado con el rol \textit{member} & RF-06 \\
    CU-03 & Propietario del equipo & Contrata un plan (pasarela en modo de pruebas) y el estado de la suscripción se refleja en la aplicación & RF-07 \\
    CU-04 & Miembro sin permiso & Intenta gestionar la facturación y se le deniega, tanto en la interfaz como en la base de datos & RF-06, RNF-02, RNF-03 \\
    CU-05 & Super-administrador & Busca un usuario, lo bloquea y lo suplanta para darle soporte & RF-08 \\
    CU-06 & Gestor de contenidos & Entra en el CMS y edita registros de una tabla según sus permisos; la acción queda auditada & RF-09, RF-10 \\
    CU-07 & Desarrollador & Arranca un nuevo SaaS a partir de PymeKit siguiendo el manual & RNF-01, RNF-07 \\
    \bottomrule
  \end{tabularx}
\end{table}

\begin{figure}[htbp]
  \centering
  \begin{tikzpicture}[
      font=\footnotesize,
      uc/.style={ellipse, draw, minimum width=4.2cm, minimum height=0.85cm,
                 align=center, inner sep=1pt},
      actor/.style={draw=none, align=center, font=\footnotesize\bfseries},
      link/.style={-}
    ]
    % Casos de uso (dentro del sistema)
    \node[uc] (cu1) at (0, 3.6)  {CU-01 Registrarse y\\acceder al área privada};
    \node[uc] (cu2) at (0, 2.4)  {CU-02 Crear equipo\\e invitar a un miembro};
    \node[uc] (cu3) at (0, 1.2)  {CU-03 Contratar\\un plan};
    \node[uc] (cu4) at (0, 0)    {CU-04 Gestionar la\\facturación (denegado)};
    \node[uc] (cu5) at (0, -1.2) {CU-05 Bloquear y\\suplantar a un usuario};
    \node[uc] (cu6) at (0, -2.4) {CU-06 Editar registros\\en el CMS (auditado)};
    \node[uc] (cu7) at (0, -3.6) {CU-07 Arrancar un\\SaaS nuevo};

    \begin{scope}[on background layer]
      \node[draw, rounded corners, fit=(cu1)(cu7), inner sep=10pt,
            label={[font=\small\bfseries]above:PymeKit}] (sys) {};
    \end{scope}

    % Actores (izquierda: clientes; derecha: personal y externos)
    \node[actor] (vis) at (-5.2, 3.6)  {Visitante};
    \node[actor] (usu) at (-5.2, 2.4)  {Usuario};
    \node[actor] (pro) at (-5.2, 1.2)  {Propietario\\del equipo};
    \node[actor] (mie) at (-5.2, 0)    {Miembro\\sin permiso};
    \node[actor] (sa)  at (5.2, -1.2)  {Super-\\administrador};
    \node[actor] (ges) at (5.2, -2.4)  {Gestor de\\contenidos};
    \node[actor] (dev) at (5.2, -3.6)  {Desarrollador};
    \node[actor] (pas) at (5.2, 1.2)   {Pasarela\\de pago};

    \draw[link] (vis) -- (cu1);
    \draw[link] (usu) -- (cu2);
    \draw[link] (pro) -- (cu3);
    \draw[link] (mie) -- (cu4);
    \draw[link] (sa)  -- (cu5);
    \draw[link] (ges) -- (cu6);
    \draw[link] (dev) -- (cu7);
    \draw[link] (pas) -- (cu3);
  \end{tikzpicture}
  \caption{Diagrama de casos de uso principales de PymeKit.}
  \label{fig:casos-uso}
  % Texto alternativo: diagrama con el sistema PymeKit y siete casos de uso
  % (CU-01 a CU-07). A la izquierda, los actores cliente (visitante, usuario,
  % propietario del equipo, miembro sin permiso) enlazados con CU-01 a CU-04;
  % a la derecha, el personal (super-administrador con CU-05, gestor de
  % contenidos con CU-06), el desarrollador con CU-07 y la pasarela de pago
  % como actor secundario de CU-03.
\end{figure}
% PENDIENTE: si la plantilla prefiere figuras de actor UML (monigote), usar la
% versión PlantUML de diagramas/casos-uso.puml y exportarla a PDF/SVG.

A continuación se describe brevemente cada caso de uso.

\begin{description}
  \item[CU-01 · Registro y acceso.] El visitante se registra con correo y
        contraseña (o con enlace mágico), recibe el correo de verificación,
        lo confirma y accede al área privada, donde dispone de su cuenta
        personal. Cubre RF-01 y RF-03.
  \item[CU-02 · Equipo e invitación.] Un usuario crea una cuenta de equipo
        para su pyme, de la que pasa a ser propietario, e invita por correo a
        un empleado con el rol \textit{member}. El invitado acepta la
        invitación y queda como miembro con los permisos de su rol. Cubre
        RF-06.
  \item[CU-03 · Contratación de un plan.] El propietario elige un plan y
        completa el pago en la pasarela (en modo de pruebas). La pasarela
        notifica el resultado mediante un \textit{webhook} y la aplicación
        refleja el estado de la suscripción de la cuenta. Cubre RF-07.
  \item[CU-04 · Acceso denegado.] Un miembro sin el permiso de gestión de la
        facturación intenta acceder a ella. La interfaz no se lo ofrece y,
        aunque forzara la petición, la base de datos la rechaza por sus
        políticas de seguridad. Es el caso que materializa la seguridad en
        profundidad (RNF-02) y el aislamiento (RNF-03).
  \item[CU-05 · Soporte desde la consola.] El super-administrador, con sesión
        de doble factor, busca un usuario en la consola, lo bloquea y lo
        suplanta temporalmente para reproducir un problema. Cubre RF-08.
  \item[CU-06 · Edición en el CMS.] Un gestor de contenidos, con un rol
        limitado del CMS, abre una tabla para la que tiene permiso, edita un
        registro y la acción queda registrada en la auditoría. Las tablas sin
        permiso no le aparecen ni puede consultarlas. Cubre RF-09 y RF-10.
  \item[CU-07 · Reutilización.] Un desarrollador arranca un SaaS nuevo a
        partir de PymeKit siguiendo el manual y cambiando solo la
        configuración. Sirve para evaluar RNF-01 y RNF-07.
        % PENDIENTE (P-01): definir el escenario y las métricas concretas.
\end{description}

\todo{Cuando termine F6, añadir para cada caso de uso su flujo principal y
alternativo (plantilla de caso de uso de la ESI, si la hay) y la prueba que lo
valida. En la fecha de este borrador, CU-01 está probado a mano por el autor;
CU-02 a CU-06 dependen de F2 y F5, y CU-07 de F6.}
